\section{Problem Statement}

%szb: ez nekem túl ismétlős, lehetne kicsit lényegretörőbb a section bevezetése
Recent advances in large language models have significantly improved the
performance of automatic code translation systems, especially for common
programming patterns and widely used language pairs. However, while these models
often produce syntactically correct and readable code, their functional
correctness is not always guaranteed, particularly in more complex scenarios
\cite{llm_limitations, correctness_llm}.

%szb: kicsit fura a sorrend. Az előbb a kiértékelés hangzott másodlagos topic-ként, míg most ez a central
%sztem ez így nemigen jó + az evaluation-t nem is igazán említeném itt, de biztos nem ennyire... Hiszen nem új benchmarkot mutatsz be.
A central challenge in this domain is the evaluation of generated code. Widely
used metrics such as BLEU measure token-level similarity between the generated
and reference code \cite{bleu_metric}, but they do not reflect whether the code
actually compiles or behaves correctly. CodeBLEU improves upon this by
incorporating syntactic and semantic features, but it still does not fully
capture functional correctness \cite{ren2020codebleu}.

To address this issue, execution-based evaluation methods have been introduced,
where generated programs are compiled and tested using predefined test cases.
These methods provide a more realistic measure of correctness, but are more
complex and computationally expensive to apply \cite{huang2022exeds}.

In addition to evaluation challenges, the impact of fine-tuning on code
translation performance is not fully understood. While fine-tuning is commonly
used to adapt pre-trained models to specific tasks, its effectiveness depends on
multiple factors, including the dataset, model architecture, and evaluation
method. %szb: nem igaz, az effectiveness nem múlik a kiértékelési method-on. Helyette az igaz, hogy: az evaluation methdo becsli az effectiveness-t
Therefore, systematic comparison between base and fine-tuned models is
necessary. %szb: esetleg mások csináltak ilyesmit? Miben több/jobb/más a te munkád, mint másoké esetleg? Lehet említeni konkrétumokat
% pl. sokan nem használnak teszt-alapú kiértékelőt, pl. sokan nem finetuneolnak, pl. sokan más adathalmazon finetuneolnak, ilyenek.

This thesis focuses on the comparative evaluation of code translation models. It
investigates how fine-tuning affects model performance, and compares different
models using both similarity-based metrics (BLEU) and execution-based evaluation
methods such as CodeEvalBench. %szb: még nem jutottam a végére, de: such as nem lehet itt. mindent le kell írni, nem csak példákat. ha meg ez az egyedüli ilyen evaluation, akkor a such as használata helytelen
The goal %szb: lehetne ink. úgy fogalmazni, hogy "az eddigieken felül az is érdekes (uninvestigated) kérdés, hogy ezek a mérések hogy viszonyulnak egymáshoz"
is to determine whether improvements observed in traditional metrics correspond
to actual functional correctness.